\section{The search state is a dependency graph}\label{sec:state}

\subsection{A branch contains more than a goal}

A minimal production branch should contain the following conceptual fields:
\begin{Code}
Branch := {
  requestId, environmentGeneration, branchId,
  rootProgram, rootContractProof,
  metaSnapshot, elaboratorSnapshot, declarationDelta,
  pendingObligations, dependencyGraph, blockedConstraints,
  recursionCapabilities, exampleBankVersion,
  domainFacts, decisionTrail, ruleCursors,
  accumulatedCost, resourceCharges
}
\end{Code}
This is a schema, not an assertion that these are stable Lean API record names.
The snapshots must include all state that the chosen rule can mutate. A purely
\MetaM{} rule can use a metaprogramming snapshot; a command-level rule creating
auxiliary declarations needs an environment/elaborator transaction as well.
Restoring only the current list of goals is insufficient. Lean's metavariable
operations and unification can assign information that later operations
observe~\cite{meta-book,lean-apply}.

Each obligation has a kind, a target expression, its owning local context, a
read/write set of reachable metavariables, dependencies, and an origin. Kinds
include computational value, type or universe choice, instance evidence,
index equality, contract proof, termination proof, and reconstruction check.
The distinction guides scheduling; it does not create different logical
notions of truth.

Connect two obligations when either can assign a metavariable reachable from
the other, including metavariables in local declaration types, local values,
instance arguments, or universe expressions. Transitive connected components
are \emph{coupled clusters}. Assignment can create or remove dependencies, so
the graph must be incrementally refreshed. A static syntactic list of initial
free metavariables is not enough.

\subsection{The fundamental continuation invariant}

Suppose a constructor creates holes $h_1:A$ and $h_2:P(h_1)$. It is incorrect to
run a solver for $h_1$, retain its first inhabitant, and only then begin an
independent solver for $h_2$. Failure of $P(h_1)$ must be able to revisit the
choice of $h_1$. The same applies to a provider's type parameter, a class
instance, an inductive index, or a recursive helper's invariant.

\begin{invariant}[Continuation ownership]\label{inv:continuation}
An alternative is successful only if its entire coupled continuation is
successful. Every decision that can affect a remaining obligation retains a
backtracking point until the cluster is accepted or deliberately suspended.
\end{invariant}

For example, the subtype
\[
 \{f:\mathbb N\to\mathbb N\to\mathbb N\mid f(11,29)=29\}
\]
first asks for a function and then for a proof. Returning the first argument is
type correct. Its proof obligation fails. Returning the second argument is then
tried and its equality reduces to reflexivity. The included Lean experiment
implements precisely this continuation discipline. It does not use a special
rule recognizing the words ``choose second.''

\subsection{Transactional rule execution}

A robust rule application follows this protocol:
\begin{Code}
tryRule(branch, rule):
  checkpoint := snapshotAllStateWritableBy(rule)
  result := rule.nextAlternative(branch)
  if result is an ordinary inapplicability:
      restore checkpoint; return NotApplicable
  if result is interruption or exhausted global resources:
      restore checkpoint; propagate status
  if result is a new typed branch:
      compute new dependencies and branch-local obligations
      retain checkpoint/decision for its coupled continuation
      return SuspendedBranch(result)
\end{Code}

Unexpected plugin exceptions should become diagnostics, not silently be treated
as a logical dead end. The microprototype intentionally has a much smaller
exception policy; production code should distinguish unification mismatch,
unsupported term form, local proof timeout, cancellation, internal error, and
resource exhaustion. In particular, repeated catching of heartbeat exhaustion
must not create a loop that appears to make search progress.

Parallel branches own independent snapshots and solver frames. Sharing
immutable expressions and environment data is useful; sharing a writable
metavariable context is not. Exporting a result from one branch to another
requires an abstracted template with explicit parameters and a fresh replay,
not copying a live expression containing foreign metavariable identifiers.

\section{Bidirectional construction over Lean expressions}\label{sec:bidirectional}

\subsection{Two judgments and a suspended constraint store}

The engine alternates between checking a proposed shape against an expected
type and inferring the type of a neutral expression:
\[
 \G\judges e\Leftarrow T
 \qquad\text{and}\qquad
 \G\judges n\Rightarrow T.
\]
A neutral expression is typically a local variable or permitted constant
applied to an argument spine, possibly followed by projections. These are
search judgments; the actual resulting term is still checked by Lean.

The expected type drives introductions. The head of an available term drives
eliminations. A suspended constraint store records equations that cannot yet be
resolved without a missing argument or instance. It is important not to
interpret ``blocked'' as ``false.'' A rule that only increases the number of
unconstrained metavariables should be charged a cost and postponed behind rules
that expose structure.

\subsection{Introduction rules}

For a dependent function, introduce a rigid local parameter:
\[
 \frac{\G,x:A\judges e\Leftarrow B(x)}
      {\G\judges (\lambda x.e)\Leftarrow \prod_{x:A}B(x)}.
\]
The binder retains its type, visibility, and instance status. A universe or
type parameter belonging to the original theorem is a rigid parameter; the
synthesizer is not allowed to specialize it to \code{Nat} merely because that
makes the current example executable.

For a dependent pair, create the witness hole before its dependent field:
\[
 \frac{\G\judges a\Leftarrow A\quad
       \G\judges b\Leftarrow B(a)}
      {\G\judges (a,b)\Leftarrow \sum_{x:A}B(x)}.
\]
Both holes remain coupled. A record constructor generalizes this rule: its
fields are processed in the actual constructor telescope, including inherited
fields and proof fields. Record synthesis must not assume that the displayed
field order is a list of independent goals.

For an inductive target $I\,\bar p\,\bar i$, inspect the environment's constructor
information. Freshen the constructor's universe levels, instantiate its
parameters, and unify its result with the target. The remaining fields become
obligations in that same transaction. A constructor whose index equation is
inconsistent can be rejected. A constructor whose equation is merely blocked
must remain available, or be recorded as a bounded search omission rather than
a proof of impossibility.

Functions and records should not always be expanded immediately. An existing
local function may solve the whole goal more cheaply than an eta expansion.
Keep a cheap exact-term lane before decomposing. Conversely, a cheap inhabitant
must not suppress later contract-sensitive alternatives. Introductory focusing
is an ordering discipline; a claim that a rule is irrevocably safe needs to
account for the whole contract and the chosen equivalence/cost notion.

\subsection{Application-spine search}\label{sec:spines}

For a candidate head $h$, do not enumerate complete argument tuples blindly.
Use a demand-directed spine algorithm:
\begin{Code}
applyHead(h, expectedType, branch):
  freshen all universe parameters of this occurrence of h
  inspect the exact dependent telescope of typeOf(h)
  enumerate admissible partial-application lengths lazily
  create fresh argument holes in telescope order
  constrain the instantiated result against expectedType
  propagate assignments through argument types and local instances
  discharge determined instance and proof arguments cheaply
  suspend blocked constraints; prioritize informative arguments
  yield the application and its entire coupled continuation
\end{Code}

Lean's \code{MVarId.apply} already provides much of the basic mechanism,
including result matching, argument metavariables, and instance processing.
The prototype uses its exact-unification configuration and requests
\code{newGoals := .all}, rather than relying on a UI-oriented goal order that
may hide dependent goals~\cite{lean-apply}. The production wrapper must also
account for delayed metavariables, universe constraints, and obligations
created by any later elaboration step. ``The tactic returned no goals'' is not
the final acceptance test.

A head may have a type such as
\[
 \prod_{A:\Type u}(x:A)\,
 \prod_{B:A\to\Type v}(k:\prod_{a:A}B(a)),\ B(x).
\]
There is no meaningful global split into ``all type arguments first, then all
value arguments.'' Instantiation follows the actual telescope. Solving $x$
may constrain $A$; the expected result may constrain $B$; a class obligation
may later determine an output parameter. All those assignments belong to one
branch.

Index providers by a conservative approximation to their result heads, arity,
and binder shape, but include fallback buckets for variable-headed and
reducible results. An index is allowed to return extra candidates. Omitting a
candidate for speed is a heuristic restriction that must not contaminate a
completeness claim. A small local context should be searched before the global
inventory, without becoming an absolute restriction on the latter.

\subsection{Higher-rank arguments and universes}

A polymorphic callback is simply an argument whose type is itself a dependent
function type. When the expected argument is
\code{(A : Type u) -> A -> A}, introduction creates a rigid $A$ and then a rigid
$x:A$; using $x$ completes that callback. A provider that accepts such a callback
can therefore be applied without translating it to a monotype or inventing a
new representation for rank two, three, or four.

Each occurrence of a polymorphic global receives fresh level metavariables.
The original query's universe parameters remain fixed. Two uses need not share
a fresh level assignment merely because they share a declaration name, while
correlations explicitly present in one application telescope must be preserved.
The prototype's independent-$u$ and independent-$v$ callbacks test one instance
of this distinction.

Lean's universe discipline must remain in force. Treating a polymorphic type
as an argument is not permission to identify \code{Type u} with an arbitrarily
large universe or to make \code{Type} impredicative. \code{Prop} has a distinct
impredicative role; nested quantification at \code{Type} still has its actual
sort~\cite{lean-types}. Lean is also noncumulative: a type at one universe does not automatically
inhabit every larger universe. An explicit lift is a typed construction, not
a numeric relaxation of a level constraint. No fixed list of universe levels
should appear in the semantic kernel of Leant2. A bounded list may be used in an explicitly
heuristic instantiation proposal lane.

Some type-valued holes remain genuinely underconstrained. Search them using
expected-type information, subexpressions of the local/provider types,
instantiation templates, and bounded enumeration of well-sorted type
expressions. Avoid committing every such hole to the first visible type.
Return an explicit request for a sketch annotation only as an interactive aid;
ordinary bounded search can still return an honest unknown.

\subsection{Unification is a service, not a completeness oracle}

Direct use of Lean removes the need to reimplement its definitional equality,
but does not make higher-order inference complete. The engine should classify
constraints into deterministic progress, blocked progress, speculative
assignment, and mismatch. A failed unification attempt is usually an
inapplicable rule alternative, not a theorem that no inhabitant exists.

In the first implementation, use Lean's standard unifier transactionally and
revisit alternative heads, argument orders where permitted, and explicit
instantiation hints. Later, add a bounded higher-order proposal layer:
pattern-like metavariable applications are solved first; remaining flexible
heads may be instantiated with small lambda/projection/application templates.
Every proposal returns to Lean for checking. This is preferable to building a
second, subtly incompatible conversion procedure.

Opaque constants, casts, and reducibility settings require particular care.
Use weak-head reduction to expose the structure needed by the current rule;
do not repeatedly normalize every expression in the environment. A reduction
cache key includes its transparency mode and relevant metavariable assignments.
Forcing an opaque provider to unfold is a policy decision, not a harmless
normalization step.

\subsection{Type classes and exact dictionary identity}

Class synthesis should use Lean's existing local-instance and registered-instance
semantics. Delayed obligations, output parameters, and instance priorities are
not replaced with Haskell-style class matching. Lean's instance search is a
specialized procedure with its own handling of such information~\cite{lean-instances}.
A wrapper should expose blocked obligations to the scheduler and retry them
when the required parameters become determined.

Two values of type \code{Payload Nat} need not contain the same payload. A
contract asking for the second dictionary's field cannot be satisfied by
silently resolving a different dictionary of the same class type. Once a
provider occurrence contains explicit evidence, keep that exact term attached
to the occurrence. This applies especially to operation dictionaries and their
law records: associativity of one multiplication does not certify another
multiplication merely because both have type $A\to A\to A$.

Instances are also a search-space policy. The normal mode should respect the
instances Lean would synthesize at that program point. An explicitly requested
evidence-enumeration mode may consider other legal dictionaries, but the emitted
term must name them and the contract must mention the same selections. Changes
to local instance scope invalidate relevant cached applications and proofs.

\section{Search scheduling, focusing, and cost}\label{sec:scheduling}

\subsection{A portfolio with a fair reference lane}

Use a weighted portfolio rather than one global best-first queue. Suitable
lanes are exact/local reuse, focused introductions and spines, indexed library
composition, recursor schemas, contract-guided sketches, and a restricted
exhaustive reference enumerator. Proof workers are a shared service with their
own bounded tasks; they are not an unlimited final filter.

Each lane exposes a resumable cursor. A scheduler grants a bounded work quantum,
not ``run until the next candidate appears.'' Otherwise a lane that produces
nothing can monopolize the request while its sibling appears fair only on
paper. Charge successful steps, failed matches, solver calls, duplicate
candidates, and cancellations that have consumed work. Do not refund computation
because its result was unhelpful.

Within a branch, choose an unblocked obligation with high constraint information
and few plausible alternatives. A known finite constructor choice, a nearly
resolved index equality, or a proof about a completed small expression can be
more informative than an unconstrained large function. Apply an aging rule so
that informative-first does not become starvation of a necessary hard goal.

The reference lane uses increasing cost bounds over an explicitly defined
finite branching grammar. It must retain all alternatives except those removed
by a valid equivalence or refutation argument. Beam truncation and aggressive
provider retrieval can be used in fast lanes, but cannot be described as fair
or complete. The report should say which lane found a result and which
restrictions were active.

\subsection{What completeness can reasonably mean}

\begin{proposition}[Relative bounded-search coverage]
Fix an environment, a finite branching rule grammar, a finite cost bound, strictly positive cost for every generative step, and a
terminating applicability/checking procedure for each bounded alternative.
Administrative zero-cost steps are required to terminate.
If the scheduler eventually visits every alternative within that bound and
only removes branches with sound no-completion certificates, every successful
leaf in that bounded search graph is eventually visited.
\end{proposition}
\begin{proof}
There are finitely many alternatives under the hypotheses. A successful leaf
has a finite path from the root. A sound no-completion certificate cannot remove
any prefix of that path. Fair visitation therefore reaches every edge of the
path and then the leaf.
\end{proof}

The theorem is intentionally about the declared search graph. It does not claim
that Lean's unifier enumerates every possible higher-order solution, that the
chosen recursion schemas express every total function, or that an SMT solver
terminates on every query. Increasing bounds yields a useful reference
semi-search only for alternatives that remain effectively enumerable and are
actually included. General proof search can be added as a further lane without
promoting its timeouts to negative evidence.

\subsection{Cost is multi-objective}

Use a structured cost, for example
\[
 (\text{policy violations},\ \text{semantic runtime cost},\
   \text{source complexity},\ \text{proof effort},\ \text{search cost}).
\]
The first component is an admission constraint, not a large numeric penalty
that can be outweighed by other terms. Runtime estimates may be unknown and
must identify their cost model. Source complexity counts constructor, branch,
application, and recursive-schema structure rather than line breaks. Proof
size can matter operationally but should not force the engine to choose an
asymptotically poor program simply because its first proof is shorter.

Different default policies are useful: first verified result for interactive
work; a small Pareto frontier for exploration; and bounded optimization for
code generation. A result is ``best within the explored set,'' unless an
optimality certificate says otherwise. Preserving alternatives matters both
for behavior and for optimization: the smallest fold encoding of reverse is
not necessarily the fastest implementation.

\subsection{Why not just call Aesop?}

Aesop supplies a mature extensible proof-search model with normalization,
indexed rules, and safe versus backtracking choices~\cite{aesop}. Leant2 should
reuse it for suitable proof obligations and borrow its engineering ideas.
However, a rule safe for proving that a type has \emph{some} inhabitant need not
be safe for enumerating computationally distinct inhabitants constrained by a
later obligation. The distinction in Invariant~\ref{inv:continuation} is the
reason to retain a synthesis-specific controller.

Similarly, \code{grind}, arithmetic tactics, and rewriting are valuable proof
services. None independently supplies the entire program grammar, recursive
invariants, example propagation, or end-to-end resource policy. Integration
means converting a completed tactic result into an exact proof term, not
trusting a solver's success flag~\cite{lean-grind}.
